\documentclass[11pt,a4paper]{article}
\usepackage[margin=1in]{geometry}
\usepackage[T1]{fontenc}
\usepackage[utf8]{inputenc}
\usepackage{times}
% Shared macros for PYNE arXiv paper series
% Include with: % Shared macros for PYNE arXiv paper series
% Include with: % Shared macros for PYNE arXiv paper series
% Include with: \input{../common/macros.tex}

\usepackage{amsmath,amssymb,amsfonts}
\usepackage{amsthm}
\usepackage{booktabs}
\usepackage{graphicx}
\usepackage{xcolor}
\usepackage{hyperref}
\usepackage{url}
\usepackage{algorithm}
\usepackage{algpseudocode}
\usepackage{listings}
\usepackage{tikz}
\usepackage{pgfplots}
\pgfplotsset{compat=1.17}
\usetikzlibrary{arrows.meta,positioning,fit,calc,shapes.geometric,shapes.multipart,decorations.pathreplacing,backgrounds,chains,matrix,patterns}
\usepackage{subcaption}
\usepackage{multirow}
\usepackage{tabularx}
\usepackage{enumitem}
\usepackage{microtype}
\usepackage{balance}
\usepackage{csquotes}
\usepackage{natbib}

% Colors
\definecolor{pyneorange}{HTML}{F97316}
\definecolor{pyneblue}{HTML}{2563EB}
\definecolor{pynegray}{HTML}{6B7280}
\definecolor{pynegreen}{HTML}{16A34A}
\definecolor{codebg}{HTML}{F8FAFC}
\definecolor{codekw}{HTML}{7C3AED}
\definecolor{codestr}{HTML}{B45309}
\definecolor{codecomment}{HTML}{64748B}

\hypersetup{
  colorlinks=true,
  linkcolor=pyneblue,
  citecolor=pynegreen,
  urlcolor=pyneorange,
  pdftitle={},
  pdfauthor={HOOX / PYNE Project}
}

% Code listings
\lstdefinestyle{pine}{
  basicstyle=\ttfamily\footnotesize,
  backgroundcolor=\color{codebg},
  keywordstyle=\color{codekw}\bfseries,
  stringstyle=\color{codestr},
  commentstyle=\color{codecomment}\itshape,
  breaklines=true,
  frame=single,
  rulecolor=\color{black!20},
  numbers=left,
  numberstyle=\tiny\color{pynegray},
  numbersep=6pt,
  xleftmargin=1.5em,
  framexleftmargin=1.2em,
  showstringspaces=false,
  tabsize=2,
  morekeywords={indicator,strategy,library,plot,plotshape,fill,hline,var,varip,if,else,for,while,switch,true,false,na,and,or,not,import,export,type,method,enum},
  morecomment=[l]{//},
  morestring=[b]",
}
\lstdefinestyle{python}{
  language=Python,
  basicstyle=\ttfamily\footnotesize,
  backgroundcolor=\color{codebg},
  keywordstyle=\color{codekw}\bfseries,
  stringstyle=\color{codestr},
  commentstyle=\color{codecomment}\itshape,
  breaklines=true,
  frame=single,
  rulecolor=\color{black!20},
  numbers=left,
  numberstyle=\tiny\color{pynegray},
  numbersep=6pt,
  xleftmargin=1.5em,
  framexleftmargin=1.2em,
  showstringspaces=false,
  tabsize=4,
}
\lstset{style=python}

% Theorem environments
\theoremstyle{plain}
\newtheorem{theorem}{Theorem}[section]
\newtheorem{lemma}[theorem]{Lemma}
\newtheorem{proposition}[theorem]{Proposition}
\newtheorem{corollary}[theorem]{Corollary}
\theoremstyle{definition}
\newtheorem{definition}[theorem]{Definition}
\newtheorem{example}[theorem]{Example}
\newtheorem{remark}[theorem]{Remark}

% Shortcuts
\newcommand{\pyne}{\textsc{Pyne}}
\newcommand{\pine}{Pine Script\texttrademark}
\newcommand{\tv}{TradingView\textregistered}
\newcommand{\asdl}{\textsc{Asdl}}
\newcommand{\antlr}{\textsc{Antlr}4}
\newcommand{\numba}{\textsc{Numba}}
\newcommand{\lsp}{\textsc{Lsp}}
\newcommand{\ohlcv}{\textsc{Ohlcv}}
\newcommand{\udt}{\textsc{Udt}}
\newcommand{\jit}{\textsc{Jit}}
\newcommand{\api}{\textsc{Api}}
\newcommand{\cli}{\textsc{Cli}}
% AST acronym (text mode). Keep amsmath's math asterisk \ast intact.
\DeclareRobustCommand{\Ast}{\textsc{Ast}}
% Papers in this series write \ast{} for the AST acronym in prose.
% Redefine robustly; use \mathast for a binary asterisk in formulas.
\let\mathast\ast
\DeclareRobustCommand{\ast}{\textsc{Ast}}
\newcommand{\ir}{\textsc{Ir}}
\newcommand{\ta}{\texttt{ta}}
\newcommand{\code}[1]{\texttt{#1}}
\newcommand{\file}[1]{\texttt{#1}}
\newcommand{\eg}{e.g.}
\newcommand{\ie}{i.e.}
\newcommand{\etal}{\textit{et~al.}}
\newcommand{\resp}{respectively}

% Math helpers
\newcommand{\R}{\mathbb{R}}
\newcommand{\N}{\mathbb{N}}
\newcommand{\Z}{\mathbb{Z}}
% Text-mode NA token for prose; use $\namath$ inside formulas if needed.
\DeclareRobustCommand{\na}{\textsf{na}}
\newcommand{\namath}{\mathsf{na}}
\newcommand{\series}[1]{\mathbf{#1}}
\newcommand{\baridx}{i}
\newcommand{\bars}{n}
\newcommand{\period}{p}
\newcommand{\abserr}{\varepsilon_{\mathrm{abs}}}
\newcommand{\relerr}{\varepsilon_{\mathrm{rel}}}
\newcommand{\maxerr}{\varepsilon_{\mathrm{max}}}
\newcommand{\meanerr}{\varepsilon_{\mathrm{mean}}}

% Pipeline notation — use inside math: $\pipeline$
\newcommand{\pipeline}{\mathrm{src}\!\to\!\mathrm{lex}\!\to\!\mathrm{parse}\!\to\!\mathrm{AST}\!\to\!\mathrm{eval}}

% Disclaimer footnote helper
\newcommand{\trademarknote}{%
  \pine{} and \tv{} are trademarks of TradingView, Inc.
  \pyne{} is an independent, unofficial implementation and is not affiliated with,
  authorized by, sponsored by, or endorsed by TradingView, Inc.
}

% Figure path helper (papers live one level under .papers/)
\graphicspath{{figures/}{../common/figures/}}


\usepackage{amsmath,amssymb,amsfonts}
\usepackage{amsthm}
\usepackage{booktabs}
\usepackage{graphicx}
\usepackage{xcolor}
\usepackage{hyperref}
\usepackage{url}
\usepackage{algorithm}
\usepackage{algpseudocode}
\usepackage{listings}
\usepackage{tikz}
\usepackage{pgfplots}
\pgfplotsset{compat=1.17}
\usetikzlibrary{arrows.meta,positioning,fit,calc,shapes.geometric,shapes.multipart,decorations.pathreplacing,backgrounds,chains,matrix,patterns}
\usepackage{subcaption}
\usepackage{multirow}
\usepackage{tabularx}
\usepackage{enumitem}
\usepackage{microtype}
\usepackage{balance}
\usepackage{csquotes}
\usepackage{natbib}

% Colors
\definecolor{pyneorange}{HTML}{F97316}
\definecolor{pyneblue}{HTML}{2563EB}
\definecolor{pynegray}{HTML}{6B7280}
\definecolor{pynegreen}{HTML}{16A34A}
\definecolor{codebg}{HTML}{F8FAFC}
\definecolor{codekw}{HTML}{7C3AED}
\definecolor{codestr}{HTML}{B45309}
\definecolor{codecomment}{HTML}{64748B}

\hypersetup{
  colorlinks=true,
  linkcolor=pyneblue,
  citecolor=pynegreen,
  urlcolor=pyneorange,
  pdftitle={},
  pdfauthor={HOOX / PYNE Project}
}

% Code listings
\lstdefinestyle{pine}{
  basicstyle=\ttfamily\footnotesize,
  backgroundcolor=\color{codebg},
  keywordstyle=\color{codekw}\bfseries,
  stringstyle=\color{codestr},
  commentstyle=\color{codecomment}\itshape,
  breaklines=true,
  frame=single,
  rulecolor=\color{black!20},
  numbers=left,
  numberstyle=\tiny\color{pynegray},
  numbersep=6pt,
  xleftmargin=1.5em,
  framexleftmargin=1.2em,
  showstringspaces=false,
  tabsize=2,
  morekeywords={indicator,strategy,library,plot,plotshape,fill,hline,var,varip,if,else,for,while,switch,true,false,na,and,or,not,import,export,type,method,enum},
  morecomment=[l]{//},
  morestring=[b]",
}
\lstdefinestyle{python}{
  language=Python,
  basicstyle=\ttfamily\footnotesize,
  backgroundcolor=\color{codebg},
  keywordstyle=\color{codekw}\bfseries,
  stringstyle=\color{codestr},
  commentstyle=\color{codecomment}\itshape,
  breaklines=true,
  frame=single,
  rulecolor=\color{black!20},
  numbers=left,
  numberstyle=\tiny\color{pynegray},
  numbersep=6pt,
  xleftmargin=1.5em,
  framexleftmargin=1.2em,
  showstringspaces=false,
  tabsize=4,
}
\lstset{style=python}

% Theorem environments
\theoremstyle{plain}
\newtheorem{theorem}{Theorem}[section]
\newtheorem{lemma}[theorem]{Lemma}
\newtheorem{proposition}[theorem]{Proposition}
\newtheorem{corollary}[theorem]{Corollary}
\theoremstyle{definition}
\newtheorem{definition}[theorem]{Definition}
\newtheorem{example}[theorem]{Example}
\newtheorem{remark}[theorem]{Remark}

% Shortcuts
\newcommand{\pyne}{\textsc{Pyne}}
\newcommand{\pine}{Pine Script\texttrademark}
\newcommand{\tv}{TradingView\textregistered}
\newcommand{\asdl}{\textsc{Asdl}}
\newcommand{\antlr}{\textsc{Antlr}4}
\newcommand{\numba}{\textsc{Numba}}
\newcommand{\lsp}{\textsc{Lsp}}
\newcommand{\ohlcv}{\textsc{Ohlcv}}
\newcommand{\udt}{\textsc{Udt}}
\newcommand{\jit}{\textsc{Jit}}
\newcommand{\api}{\textsc{Api}}
\newcommand{\cli}{\textsc{Cli}}
% AST acronym (text mode). Keep amsmath's math asterisk \ast intact.
\DeclareRobustCommand{\Ast}{\textsc{Ast}}
% Papers in this series write \ast{} for the AST acronym in prose.
% Redefine robustly; use \mathast for a binary asterisk in formulas.
\let\mathast\ast
\DeclareRobustCommand{\ast}{\textsc{Ast}}
\newcommand{\ir}{\textsc{Ir}}
\newcommand{\ta}{\texttt{ta}}
\newcommand{\code}[1]{\texttt{#1}}
\newcommand{\file}[1]{\texttt{#1}}
\newcommand{\eg}{e.g.}
\newcommand{\ie}{i.e.}
\newcommand{\etal}{\textit{et~al.}}
\newcommand{\resp}{respectively}

% Math helpers
\newcommand{\R}{\mathbb{R}}
\newcommand{\N}{\mathbb{N}}
\newcommand{\Z}{\mathbb{Z}}
% Text-mode NA token for prose; use $\namath$ inside formulas if needed.
\DeclareRobustCommand{\na}{\textsf{na}}
\newcommand{\namath}{\mathsf{na}}
\newcommand{\series}[1]{\mathbf{#1}}
\newcommand{\baridx}{i}
\newcommand{\bars}{n}
\newcommand{\period}{p}
\newcommand{\abserr}{\varepsilon_{\mathrm{abs}}}
\newcommand{\relerr}{\varepsilon_{\mathrm{rel}}}
\newcommand{\maxerr}{\varepsilon_{\mathrm{max}}}
\newcommand{\meanerr}{\varepsilon_{\mathrm{mean}}}

% Pipeline notation — use inside math: $\pipeline$
\newcommand{\pipeline}{\mathrm{src}\!\to\!\mathrm{lex}\!\to\!\mathrm{parse}\!\to\!\mathrm{AST}\!\to\!\mathrm{eval}}

% Disclaimer footnote helper
\newcommand{\trademarknote}{%
  \pine{} and \tv{} are trademarks of TradingView, Inc.
  \pyne{} is an independent, unofficial implementation and is not affiliated with,
  authorized by, sponsored by, or endorsed by TradingView, Inc.
}

% Figure path helper (papers live one level under .papers/)
\graphicspath{{figures/}{../common/figures/}}


\usepackage{amsmath,amssymb,amsfonts}
\usepackage{amsthm}
\usepackage{booktabs}
\usepackage{graphicx}
\usepackage{xcolor}
\usepackage{hyperref}
\usepackage{url}
\usepackage{algorithm}
\usepackage{algpseudocode}
\usepackage{listings}
\usepackage{tikz}
\usepackage{pgfplots}
\pgfplotsset{compat=1.17}
\usetikzlibrary{arrows.meta,positioning,fit,calc,shapes.geometric,shapes.multipart,decorations.pathreplacing,backgrounds,chains,matrix,patterns}
\usepackage{subcaption}
\usepackage{multirow}
\usepackage{tabularx}
\usepackage{enumitem}
\usepackage{microtype}
\usepackage{balance}
\usepackage{csquotes}
\usepackage{natbib}

% Colors
\definecolor{pyneorange}{HTML}{F97316}
\definecolor{pyneblue}{HTML}{2563EB}
\definecolor{pynegray}{HTML}{6B7280}
\definecolor{pynegreen}{HTML}{16A34A}
\definecolor{codebg}{HTML}{F8FAFC}
\definecolor{codekw}{HTML}{7C3AED}
\definecolor{codestr}{HTML}{B45309}
\definecolor{codecomment}{HTML}{64748B}

\hypersetup{
  colorlinks=true,
  linkcolor=pyneblue,
  citecolor=pynegreen,
  urlcolor=pyneorange,
  pdftitle={},
  pdfauthor={HOOX / PYNE Project}
}

% Code listings
\lstdefinestyle{pine}{
  basicstyle=\ttfamily\footnotesize,
  backgroundcolor=\color{codebg},
  keywordstyle=\color{codekw}\bfseries,
  stringstyle=\color{codestr},
  commentstyle=\color{codecomment}\itshape,
  breaklines=true,
  frame=single,
  rulecolor=\color{black!20},
  numbers=left,
  numberstyle=\tiny\color{pynegray},
  numbersep=6pt,
  xleftmargin=1.5em,
  framexleftmargin=1.2em,
  showstringspaces=false,
  tabsize=2,
  morekeywords={indicator,strategy,library,plot,plotshape,fill,hline,var,varip,if,else,for,while,switch,true,false,na,and,or,not,import,export,type,method,enum},
  morecomment=[l]{//},
  morestring=[b]",
}
\lstdefinestyle{python}{
  language=Python,
  basicstyle=\ttfamily\footnotesize,
  backgroundcolor=\color{codebg},
  keywordstyle=\color{codekw}\bfseries,
  stringstyle=\color{codestr},
  commentstyle=\color{codecomment}\itshape,
  breaklines=true,
  frame=single,
  rulecolor=\color{black!20},
  numbers=left,
  numberstyle=\tiny\color{pynegray},
  numbersep=6pt,
  xleftmargin=1.5em,
  framexleftmargin=1.2em,
  showstringspaces=false,
  tabsize=4,
}
\lstset{style=python}

% Theorem environments
\theoremstyle{plain}
\newtheorem{theorem}{Theorem}[section]
\newtheorem{lemma}[theorem]{Lemma}
\newtheorem{proposition}[theorem]{Proposition}
\newtheorem{corollary}[theorem]{Corollary}
\theoremstyle{definition}
\newtheorem{definition}[theorem]{Definition}
\newtheorem{example}[theorem]{Example}
\newtheorem{remark}[theorem]{Remark}

% Shortcuts
\newcommand{\pyne}{\textsc{Pyne}}
\newcommand{\pine}{Pine Script\texttrademark}
\newcommand{\tv}{TradingView\textregistered}
\newcommand{\asdl}{\textsc{Asdl}}
\newcommand{\antlr}{\textsc{Antlr}4}
\newcommand{\numba}{\textsc{Numba}}
\newcommand{\lsp}{\textsc{Lsp}}
\newcommand{\ohlcv}{\textsc{Ohlcv}}
\newcommand{\udt}{\textsc{Udt}}
\newcommand{\jit}{\textsc{Jit}}
\newcommand{\api}{\textsc{Api}}
\newcommand{\cli}{\textsc{Cli}}
% AST acronym (text mode). Keep amsmath's math asterisk \ast intact.
\DeclareRobustCommand{\Ast}{\textsc{Ast}}
% Papers in this series write \ast{} for the AST acronym in prose.
% Redefine robustly; use \mathast for a binary asterisk in formulas.
\let\mathast\ast
\DeclareRobustCommand{\ast}{\textsc{Ast}}
\newcommand{\ir}{\textsc{Ir}}
\newcommand{\ta}{\texttt{ta}}
\newcommand{\code}[1]{\texttt{#1}}
\newcommand{\file}[1]{\texttt{#1}}
\newcommand{\eg}{e.g.}
\newcommand{\ie}{i.e.}
\newcommand{\etal}{\textit{et~al.}}
\newcommand{\resp}{respectively}

% Math helpers
\newcommand{\R}{\mathbb{R}}
\newcommand{\N}{\mathbb{N}}
\newcommand{\Z}{\mathbb{Z}}
% Text-mode NA token for prose; use $\namath$ inside formulas if needed.
\DeclareRobustCommand{\na}{\textsf{na}}
\newcommand{\namath}{\mathsf{na}}
\newcommand{\series}[1]{\mathbf{#1}}
\newcommand{\baridx}{i}
\newcommand{\bars}{n}
\newcommand{\period}{p}
\newcommand{\abserr}{\varepsilon_{\mathrm{abs}}}
\newcommand{\relerr}{\varepsilon_{\mathrm{rel}}}
\newcommand{\maxerr}{\varepsilon_{\mathrm{max}}}
\newcommand{\meanerr}{\varepsilon_{\mathrm{mean}}}

% Pipeline notation — use inside math: $\pipeline$
\newcommand{\pipeline}{\mathrm{src}\!\to\!\mathrm{lex}\!\to\!\mathrm{parse}\!\to\!\mathrm{AST}\!\to\!\mathrm{eval}}

% Disclaimer footnote helper
\newcommand{\trademarknote}{%
  \pine{} and \tv{} are trademarks of TradingView, Inc.
  \pyne{} is an independent, unofficial implementation and is not affiliated with,
  authorized by, sponsored by, or endorsed by TradingView, Inc.
}

% Figure path helper (papers live one level under .papers/)
\graphicspath{{figures/}{../common/figures/}}


\title{%
  \textbf{PYNE: An Open Dual-Engine Toolchain\\for Offline Pine Script Evaluation}%
}
\author{%
  jango\_blockchained\thanks{Correspondence: \texttt{contact@hoox.sh}}
  \and
  HOOX Project contributors\\
  \small Independent Open Source Research\\
  \small \url{https://hoox.sh}~\textperiodcentered~\url{https://github.com/hoox-sh/pyne}
}
\date{2026}

\begin{document}
\maketitle

\begin{abstract}
Domain-specific languages (\textsc{Dsl}s) for technical analysis are often
tightly coupled to proprietary charting platforms, which limits reproducible
research, continuous integration, and editor-centric development workflows.
We present \pyne{} (PyPI package \texttt{hoox-pyne}, import name
\texttt{pynescript}), an independent, unofficial open toolchain that models
\pine{} (v5--v6) as an inspectable offline pipeline:
source text is lexed and parsed with \antlr{}, lowered to an
\asdl{} abstract syntax tree (\ast{}), and evaluated under a deterministic
bar-loop runtime with a dual execution engine---an \ast{} interpreter and a
\numba{}-backed compiler with nopython kernels and object-mode fallback.
The same \ast{} underpins a desk \cli{}, a Language Server Protocol (\lsp{})
implementation with hundreds of builtins, a Flask Pro \api{}, editor
integrations, and edge/browser workers.
We formalize series semantics, \na{} propagation, and strategy fill behaviour
as implemented by the runtime; describe warm-compile caching and security
policy for \code{request.security}; and report internal performance
measurements (2026) showing order-of-magnitude gains for numeric scripts under
the compile path relative to interpretation.
\pyne{} is licensed under \textsc{Agpl}-3.0-or-later and is part of the
\textsc{Hoox} open trading stack. It is not affiliated with TradingView.
\footnote{\trademarknote}
\end{abstract}

\noindent\textbf{Keywords:}
domain-specific languages, financial time series, dual-engine runtimes,
language toolchains, Pine Script, open source compilers

\tableofcontents
\newpage

%==============================================================================
\section{Introduction}
\label{sec:intro}
%==============================================================================

\subsection{Motivation}
\label{sec:motivation}

\pine{} is a widely used domain-specific language for indicators and trading
strategies~\citep{tradingview-pine,murphy1999technical}.
In common practice, scripts are authored and executed inside a browser-locked
host environment: the editor, chart surface, historical data plane, and
execution engine are provided as a single product surface.
While this coupling optimizes interactive charting, it creates friction for
several engineering activities that are standard elsewhere in software:

\begin{itemize}[leftmargin=1.2em]
  \item \textbf{Reproducible offline evaluation.}
    Batch backtests, property-based tests, and continuous integration require
    deterministic execution without a live UI session.
  \item \textbf{Editor-native tooling.}
    Language servers~\citep{lsp2023,voutilainen2016lsp} expect parse trees,
    diagnostics, and symbol tables in-process or over a local protocol, not
    only inside a remote web IDE.
  \item \textbf{Program analysis and transformation.}
    Static linting, formatting-as-unparsing, and compile-time lowering benefit
    from a first-class \ast{} rather than opaque host evaluation.
  \item \textbf{Heterogeneous deployment.}
    The same evaluation contract is useful at the desk (\cli{}), on a
    server (\api{}), at the edge (workers), and in the browser
    (\textsc{Pyodide}/\textsc{Axis}).
\end{itemize}

\pyne{} addresses these needs by treating \pine{} as a language with a
classical compiler pipeline~\citep{aho2006dragon,cooper2011engineering,nystrom2021crafting}
rather than solely as a feature of a charting application.
The project is an \emph{independent, unofficial} implementation: it aims for
interoperability with public language documentation and community scripts, not
for binary compatibility with proprietary platform internals~\citep{pynescript2026,hoox2026}.

\subsection{Contributions}
\label{sec:contributions}

This paper makes the following contributions:

\begin{enumerate}[leftmargin=1.4em]
  \item \textbf{System architecture.}
    A layered architecture in which editors, \lsp{}, core package, dual-engine
    runtime, Pro \api{}, and edge workers share one \ast{} and one evaluate
    contract (Section~\ref{sec:architecture}).
  \item \textbf{Language front-end.}
    An \antlr{} grammar approximating \pine{} v5--v6, \asdl{}-generated nodes,
    visitor/transformer patterns, and lossless parse$\leftrightarrow$unparse
    round-trip (Section~\ref{sec:frontend}; grammar and series details are
    expanded in sister papers of this series).
  \item \textbf{Dual-engine runtime.}
    A bar-loop host with modes
    $\in \{\texttt{auto},\texttt{compile},\texttt{interpret}\}$, warm \ir{}
    cache, object-mode fallback, and a uniform response envelope for plots,
    drawings, strategy events, and alerts (Section~\ref{sec:runtime}).
  \item \textbf{Product surfaces.}
    \cli{}, \lsp{} ($\sim$800+ builtins), VS~Code extension, Flask Pro \api{},
    and editor/edge clients built on the shared core (Section~\ref{sec:surfaces}).
  \item \textbf{Evaluation experience.}
    Internal benchmarks (2026) for parse, unparse, interpret, and compile paths;
    compatibility posture; and deployment lessons (Section~\ref{sec:eval}).
\end{enumerate}

\subsection{Paper roadmap}
\label{sec:roadmap}

Section~\ref{sec:background} reviews the \pine{} bar model and series
semantics.
Section~\ref{sec:architecture} presents the end-to-end system.
Sections~\ref{sec:frontend}--\ref{sec:surfaces} detail the front-end, runtime,
and product surfaces.
Section~\ref{sec:eval} reports measurements and operational experience.
Section~\ref{sec:related} situates \pyne{} among \textsc{Dsl} toolchains and
financial languages.
Section~\ref{sec:discussion} discusses limitations and licensing implications.
Section~\ref{sec:conclusion} concludes.

%==============================================================================
\section{Background}
\label{sec:background}
%==============================================================================

\subsection{Bar model and execution cadence}
\label{sec:bar-model}

\pine{} programs are not general-purpose stream processors with arbitrary
event loops.
They are evaluated against a finite sequence of market bars.
Let a bar series of length $n$ be
\begin{equation}
  B = \bigl(b_0, b_1, \ldots, b_{n-1}\bigr),
  \qquad
  b_{\baridx} = \bigl(o_{\baridx}, h_{\baridx}, l_{\baridx}, c_{\baridx}, v_{\baridx}, t_{\baridx}\bigr),
  \label{eq:bars}
\end{equation}
where $o,h,l,c,v,t$ denote open, high, low, close, volume, and timestamp
(\ohlcv{} plus time).
The integer \emph{bar index} $\baridx \in \{0,\ldots,n-1\}$ is exposed to
scripts as \code{bar\_index}.

\begin{definition}[Bar-loop evaluation]
\label{def:barloop}
Given a script $P$ and bar sequence $B$, bar-loop evaluation produces a run
trace $\mathcal{T}$ by invoking the script body once for each
$\baridx = 0,\ldots,n-1$ in increasing order, with built-in series
(\code{open}, \code{high}, \code{low}, \code{close}, \code{volume}, \ldots)
bound to the corresponding columns of $B$ at the current index, and with
mutable series storage updated after each bar.
\end{definition}

This model matches classical technical analysis practice, in which indicators
are defined as causal functions of historical prices~\citep{wilder1978new,box1970time,tsay2010analysis}.

\subsection{Series semantics}
\label{sec:series}

Many values in \pine{} are \emph{series}: sequences indexed by bar history.

\begin{definition}[Series]
\label{def:series}
A series of type $\tau$ is a partial function
$s: \N_0 \rightharpoonup \tau \cup \{\na\}$ where $s(k)$ denotes the value
$k$ bars ago relative to the current bar index ($k=0$ is the current bar).
Writing $s[k]$ in source notation denotes historical indexing.
\end{definition}

On each bar, the runtime \emph{pushes} a current value onto each live series
and retains a finite history window.
\pyne{} optionally caps history length (series cap) and can store history in a
ring buffer to bound memory on long runs (Section~\ref{sec:eval}).

\subsection{The $\mathsf{na}$ value and propagation}
\label{sec:na}

\pine{} uses a dedicated missingness token $\na$ (``not available''), distinct
from IEEE floating-point NaN though related in spirit~\citep{ieee754,higham2002accuracy}.

\begin{definition}[$\mathsf{na}$-propagation]
\label{def:na}
For arithmetic operators $\oplus \in \{+,-,\times,/\}$ and comparisons
$\bowtie$, if either operand is $\na$ then the result is $\na$
(with documented exceptions for short-circuit Boolean operators).
Built-ins that require a full lookback window of valid samples typically
return $\na$ until sufficient history exists.
\end{definition}

Short-circuit \code{and}/\code{or} evaluate only as many operands as needed;
this interacts with side-effecting expressions and with $\na$ in Boolean
contexts. \pyne{} implements these rules in both engines and validates them
with a parity harness (Section~\ref{sec:runtime}).

\subsection{Persistence: \texttt{var} and \texttt{varip}}
\label{sec:var}

Two declaration forms control cross-bar state:

\begin{itemize}[leftmargin=1.2em]
  \item \code{var} initializes on the first bar (conceptually when
    $\baridx = 0$) and retains the value across subsequent bars unless
    reassigned.
  \item \code{varip} retains values across realtime tick updates within a
    forming bar in host environments that distinguish tick vs.\ bar-close
    evaluation; offline historical runs treat \code{varip} analogously to
    retained state with host-specific tick semantics when ticks are supplied.
\end{itemize}

These constructs are essential for accumulators, state machines, and strategy
position bookkeeping.

\subsection{Indicators versus strategies}
\label{sec:ind-vs-strat}

An \emph{indicator} script primarily produces visual and analytic series:
\code{plot}, \code{plotshape}, \code{fill}, \code{hline}, drawings, and
alerts.
A \emph{strategy} script additionally emits order intents
(\code{strategy.entry}, \code{strategy.close}, \ldots) that a broker
simulation fills subject to commission, slippage, pyramiding, and pending-fill
rules.
\pyne{} implements both surfaces under one runtime host, with strategy events
exported in the response envelope (Section~\ref{sec:runtime}).

\subsection{Illustrative program}
\label{sec:example}

Listing~\ref{lst:rsi} shows a minimal indicator.
Under bar-loop evaluation, \code{ta.rsi} maintains incremental Wilder-style
state~\citep{wilder1978new}, and \code{plot} records one output sample per bar.

\begin{lstlisting}[style=pine,caption={Minimal RSI indicator (Pine).},label={lst:rsi}]
//@version=6
indicator("RSI Demo", overlay=false)
len = input.int(14, "Length")
r = ta.rsi(close, len)
plot(r, "RSI")
hline(70)
hline(30)
\end{lstlisting}

%==============================================================================
\section{System Architecture}
\label{sec:architecture}
%==============================================================================

\subsection{Design goals}
\label{sec:goals}

The architecture is driven by four goals:

\begin{enumerate}[leftmargin=1.4em]
  \item \textbf{Single semantic core.}
    Parsing, static checks, interpretation, and compilation share one \ast{}
    schema so that editor diagnostics and Pro evaluation cannot diverge in
    structure.
  \item \textbf{Engine substitutability.}
    Clients select \texttt{interpret}, \texttt{compile}, or \texttt{auto}
    without changing the external response schema.
  \item \textbf{Small hand-edited surface.}
    Only grammar resources (\texttt{*.g4}) and the \asdl{} schema are
    hand-maintained; lexer/parser and node classes are generated~\citep{parr2013antlr,wang1997asdl}.
  \item \textbf{Deployability.}
    The same package installs from PyPI, runs as \cli{}/\lsp{} binaries, and
    powers HTTP and edge deployments~\citep{hoox2026}.
\end{enumerate}

\subsection{Layered overview}
\label{sec:layers}

Figure~\ref{fig:architecture} summarizes the system.
Editors communicate with \texttt{pynescript-lsp} (a pygls-based
server~\citep{murphy2020pygls,lsp2023}) over STDIO.
The server parses source into the shared \ast{}, runs the linter, and
services completion, hover, formatting (via unparse), definition/references,
and document symbols.
The core package (\texttt{pynescript}) owns the grammar, \ast{} builder,
evaluators, and compiler.
The dual-engine runtime executes scripts against \ohlcv{} data and emits
plots, drawings, strategy events, and alerts.
The Pro \api{} (Flask) exposes \texttt{POST /run}, chart preview, and quick
backtest endpoints on the same evaluation path.
Edge workers and browser runtimes (Pyodide/\textsc{Axis}) reuse the evaluate
contract with thinner hosts.

\begin{figure}[t]
  \centering
  \resizebox{\textwidth}{!}{% End-to-end PYNE system architecture (TikZ)
\begin{tikzpicture}[
  font=\sffamily\footnotesize,
  box/.style={draw, rounded corners=2pt, align=center, inner sep=5pt, minimum height=1.1cm},
  layer/.style={box, fill=#1!12, draw=#1!70!black, thick},
  arr/.style={-{Stealth[length=2.2mm]}, thick, pynegray},
  lbl/.style={font=\sffamily\tiny\color{pynegray}},
  node distance=0.55cm and 0.7cm
]
  % Editors
  \node[layer=pyneblue, minimum width=2.4cm] (vscode) {VS Code\\ext.};
  \node[layer=pyneblue, minimum width=2.0cm, right=of vscode] (nvim) {Neovim};
  \node[layer=pyneblue, minimum width=1.8cm, right=of nvim] (zed) {Zed};
  \node[layer=pyneblue, minimum width=1.8cm, right=of zed] (emacs) {Emacs};
  \node[layer=pyneblue, minimum width=2.2cm, right=of emacs] (cli) {\cli{} / SDK};

  \node[fit=(vscode)(nvim)(zed)(emacs)(cli), draw=pyneblue!40, dashed, inner sep=6pt,
        label={[lbl]above:Editors \& clients}] (editors) {};

  % LSP
  \node[layer=pyneorange, below=1.0cm of editors, minimum width=11.5cm, minimum height=1.3cm] (lsp) {%
    \textbf{\texttt{pynescript-lsp}} (pygls)\\[1pt]
    Workspace $\to$ Parser $\to$ \ast{} $\to$ Linter / Completion / Hover / Format / Nav
  };
  \draw[arr] (editors.south) -- node[right,lbl]{LSP over STDIO} (lsp.north);

  % Core package
  \node[layer=pynegreen, below=0.9cm of lsp, minimum width=5.4cm, minimum height=2.0cm, xshift=-3.1cm] (core) {%
    \textbf{Core package}\\\texttt{pynescript}\\[2pt]
    ANTLR4 $\to$ ASDL \ast{}\\
    Visitor / Transformer\\
    Linter + type system
  };
  \node[layer=pyneorange, below=0.9cm of lsp, minimum width=5.4cm, minimum height=2.0cm, xshift=3.1cm] (runtime) {%
    \textbf{Dual-engine runtime}\\[2pt]
    interpret $\mid$ compile $\mid$ auto\\
    Bar-loop + series store\\
    Strategy / plots / alerts
  };
  \draw[arr] (lsp.south) -- ++(0,-0.25) -| (core.north);
  \draw[arr] (lsp.south) -- ++(0,-0.25) -| (runtime.north);
  \draw[arr] (core.east) -- node[above,lbl]{shared \ast{}} (runtime.west);

  % Grammar artifacts
  \node[layer=pynegray, below=0.7cm of core, minimum width=5.4cm] (gram) {%
    Hand-edited: \texttt{*.g4}, \texttt{Pinescript.asdl}\\
    Generated: lexer/parser, ASDL nodes
  };
  \draw[arr] (gram.north) -- (core.south);

  % Pro API
  \node[layer=pyneblue, below=0.7cm of runtime, minimum width=5.4cm, minimum height=1.6cm] (api) {%
    \textbf{Pro \api{} (Flask)}\\
    \texttt{POST /run}, \texttt{/preview/chart},\\
    \texttt{/backtest/quick}
  };
  \draw[arr] (runtime.south) -- (api.north);

  % Workers / edge
  \node[layer=pynegreen, below=0.85cm of gram, minimum width=5.4cm, xshift=0cm] (edge) {%
    Edge / workers\\
    pine-worker · pyne-worker\\
    Pyodide / AXIS (browser)
  };
  \draw[arr] (core.south west) ++(0.3,-0.05) |- (edge.north west);
  \draw[arr] (runtime.south east) ++(-0.3,-0.05) |- (edge.north east);
  \draw[arr] (api.south) -- ++(0,-0.35) -| (edge.east);

  % Outputs
  \node[layer=pyneorange, below=0.7cm of api, minimum width=5.4cm, xshift=0cm] (out) {%
    Plots · fills · drawings\\
    Strategy events · alerts\\
    Chart PNG / equity JSON
  };
  \draw[arr] (api.south) -- (out.north);
\end{tikzpicture}
}
  \caption{End-to-end \pyne{} architecture: editors and clients, \lsp{},
  core package and dual-engine runtime, Pro \api{}, and edge/browser workers.
  Hand-edited artifacts are the ANTLR grammars and ASDL schema; generated
  code includes the lexer/parser and ASDL node classes.}
  \label{fig:architecture}
\end{figure}

\subsection{Pipeline notation}
\label{sec:pipeline}

The canonical offline pipeline is
\begin{equation}
  \text{source}
  \xrightarrow{\text{lex/parse}}
  \text{parse tree}
  \xrightarrow{\text{builder}}
  \text{\ast{}}
  \xrightarrow{\text{bar-loop}}
  \text{artifacts},
  \label{eq:pipeline}
\end{equation}
where \emph{artifacts} denote plots, fills, drawings, strategy events, and
alerts.
Source suffixes \texttt{.pine} and \texttt{.pyne} are accepted equivalently.
Optionally, clients post-process artifacts (PNG charts, equity curves, webhook
alerts).

\subsection{Hand-edited versus generated code}
\label{sec:hand-gen}

Following classic compiler engineering practice~\citep{aho2006dragon,appeli1997modern},
\pyne{} minimizes the trusted hand-written front-end surface:

\begin{itemize}[leftmargin=1.2em]
  \item \textbf{Hand-edited:}
    \file{antlr4/resource/*.g4},
    \file{asdl/resource/Pinescript.asdl},
    builder/visitor/evaluator/compiler logic.
  \item \textbf{Generated:}
    Python lexer and parser under \file{antlr4/generated/},
    ASDL node classes under \file{asdl/generated/}.
\end{itemize}

Committed generated artifacts allow end users to install from PyPI without a
Java toolchain, while maintainers regenerate from resources when the language
surface evolves (e.g., v6 multiline strings).

\subsection{Shared AST between editor and Pro API}
\label{sec:shared-ast}

A recurring failure mode in dual-product systems is semantic drift between
``IDE parsing'' and ``server execution.''
\pyne{} avoids this by construction: the Pro \api{} and the \lsp{} both call
into the same \texttt{pynescript} package.
The Pro path adds authentication, data binding, chart rendering, and
operational controls (timeouts, prewarm), not a second language
implementation~\citep{flask}.

%==============================================================================
\section{Language Front-End}
\label{sec:frontend}
%==============================================================================

This section gives a high-level account of the front-end.
Sister papers in the series expand grammar engineering, series storage, and
parity methodology; here we emphasize interfaces required by the architecture.

\subsection{Grammar}
\label{sec:grammar}

The lexer and parser are specified in \antlr{}~4~\citep{parr2013antlr,parr2011allstar,antlr4}.
The grammar approximates the public \pine{} v5--v6 surface, including:

\begin{itemize}[leftmargin=1.2em]
  \item Script kinds: \code{indicator}, \code{strategy}, \code{library}.
  \item Declarations: functions, \code{type} user-defined types (\udt{}),
    \code{enum}, \code{method}, imports/exports.
  \item Control flow: \code{if}/\code{else}, \code{for}/\code{for\ldots to},
    \code{while}, \code{switch}, \code{break}/\code{continue}.
  \item Literals: numbers, strings (including v6 triple-quoted multiline),
    colours, \code{na}, booleans.
  \item Significant indentation interactions handled via a hand-written lexer
    base class cooperating with generated rules.
\end{itemize}

Practical grammar evolution (e.g., triple-quoted strings) requires careful
fragment design, selective regeneration of lexer artifacts, and round-trip
fixtures~\citep{pynescript2026}.

\subsection{ASDL AST}
\label{sec:asdl}

Node types are defined in an \asdl{} schema~\citep{wang1997asdl,asdl-python}
and generated into Python classes.
The root \code{Script} node carries a statement list and annotation strings
(e.g., \code{//@version=6}) promoted from comments so that parse$\to$unparse
preserves version and description metadata.

\begin{lstlisting}[style=python,caption={Parse and unparse round-trip (Python API).},label={lst:roundtrip}]
from pynescript.ast.helper import parse, unparse

src = '''//@version=6
indicator("Demo")
plot(close)
'''
tree = parse(src)
assert "plot" in unparse(tree)
\end{lstlisting}

Visitors and transformers~\citep{gamma1994design,gamma-visitor} provide the
extension points used by the linter, compiler emitter, and evaluators.

\subsection{Round-trip fidelity}
\label{sec:roundtrip}

Lossless round-trip is a project invariant for supported syntax:
\begin{equation}
  \mathrm{unparse}(\mathrm{parse}(s)) \equiv_{\mathrm{syn}} s',
  \label{eq:roundtrip}
\end{equation}
where $\equiv_{\mathrm{syn}}$ denotes syntactic equivalence up to documented
normalization (whitespace policies, annotation attachment).
Round-trip underpins formatting-as-unparsing in the \lsp{} and guards against
silent \ast{} incompleteness when new syntax is added.

\subsection{Linter and type system}
\label{sec:linter}

The linter implements a set of structural and style rules (undeclared names,
suspicious assignments, version issues, etc.) on the \ast{}.
A type system layer tracks primitive types, series wrappers, and \udt{} field
shapes for completion and diagnostics~\citep{pierce2002types,cardelli1996type}.
Full formalization of the type rules is out of scope for this architecture
paper; the operational interface is: parse once, lint and complete from the
same tree, evaluate from the same tree.

\subsection{Builtin metadata}
\label{sec:builtins}

Builtin completion and hover draw from structured metadata covering namespaces
such as \ta{}, \code{math}, \code{str}, \code{array}, \code{matrix},
\code{map}, \code{strategy}, \code{request}, and drawing objects.
The \lsp{} exposes on the order of $800+$ completion items, exceeding early
inventory counts as v6 surface area grew~\citep{pynescript2026}.

%==============================================================================
\section{Dual-Engine Runtime}
\label{sec:runtime}
%==============================================================================

\subsection{Execution modes}
\label{sec:modes}

The runtime accepts an explicit mode parameter
\begin{equation}
  \textit{mode} \in \{\texttt{auto},\, \texttt{compile},\, \texttt{interpret}\}.
  \label{eq:modes}
\end{equation}

\begin{itemize}[leftmargin=1.2em]
  \item \textbf{Interpret.}
    A classical \ast{} walk~\citep{nystrom2021crafting}: each bar re-enters
    the script body; builtins dispatch through Python mixins; series are
    updated in the host.
  \item \textbf{Compile.}
    A compiler visitor lowers supported numeric subgraphs to
    \numba{}~\citep{lam2015numba,lattner2004llvm} nopython kernels (and
    related helpers), with object-mode regions for features that are not yet
    nopython-safe (strategy bookkeeping, some drawings, dynamic Python
    objects).
  \item \textbf{Auto.}
    Prefer compile when feasible; on hard failure or unsupported lowering,
    fall back to interpret while preserving the response envelope.
\end{itemize}

Figure~\ref{fig:dual} shows the decision flow for warm cache, nopython
success, object-mode fallback, and interpret recovery.

\begin{figure}[t]
  \centering
  \resizebox{0.98\textwidth}{!}{% Dual-engine decision flowchart (auto mode + fallback)
\begin{tikzpicture}[
  font=\sffamily\footnotesize,
  box/.style={draw, rounded corners=2pt, align=center, inner sep=5pt, minimum width=2.8cm},
  dec/.style={diamond, draw, aspect=2.2, align=center, inner sep=1pt, fill=pyneorange!12, draw=pyneorange!70!black},
  proc/.style={box, fill=pyneblue!10, draw=pyneblue!70!black},
  ok/.style={box, fill=pynegreen!12, draw=pynegreen!70!black},
  fail/.style={box, fill=red!8, draw=red!50!black},
  arr/.style={-{Stealth[length=2mm]}, thick},
  node distance=0.55cm and 0.9cm
]
  \node[proc] (src) {Source + options\\(\texttt{mode}, OHLCV)};
  \node[dec, below=of src] (mode) {\texttt{mode}?};
  \node[proc, below left=0.9cm and 1.6cm of mode] (interp) {Interpret\\AST walk};
  \node[proc, below right=0.9cm and 1.6cm of mode] (compile) {Compile path};
  \node[dec, below=0.7cm of compile] (warm) {Warm cache\\hit?};
  \node[proc, right=1.1cm of warm] (jit) {Numba JIT\\nopython emit};
  \node[proc, below=0.65cm of warm] (load) {Load disk/\allowbreak process IR};
  \node[dec, below=0.7cm of load] (nopy) {nopython\\OK?};
  \node[proc, right=1.1cm of nopy] (obj) {Object-mode\\fallback};
  \node[ok, below=0.7cm of nopy] (exec) {Execute kernels\\bar-loop};
  \node[ok, below=1.55cm of interp] (env) {Response envelope\\plots / events / meta};
  \node[fail, below=0.55cm of obj] (fb) {On hard failure:\\fall back to interpret};
  \node[proc, left=of mode, xshift=-0.2cm] (auto) {\texttt{auto}\\prefer compile};

  \draw[arr] (src) -- (mode);
  \draw[arr] (mode) -- node[above left, font=\tiny]{\texttt{interpret}} (interp);
  \draw[arr] (mode) -- node[above right, font=\tiny]{\texttt{compile}} (compile);
  \draw[arr] (mode.west) -- node[above, font=\tiny]{\texttt{auto}} (auto.east);
  \draw[arr] (auto) |- (compile);
  \draw[arr] (compile) -- (warm);
  \draw[arr] (warm) -- node[above, font=\tiny]{yes} (load);
  \draw[arr] (warm) -- node[above, font=\tiny]{no} (jit);
  \draw[arr] (jit) |- (load);
  \draw[arr] (load) -- (nopy);
  \draw[arr] (nopy) -- node[left, font=\tiny]{yes} (exec);
  \draw[arr] (nopy) -- node[above, font=\tiny]{no} (obj);
  \draw[arr] (obj) -- (fb);
  \draw[arr] (fb.west) -| node[near start, left, font=\tiny]{} (interp.south);
  \draw[arr] (exec.west) -| (env.north);
  \draw[arr] (interp) -- (env);
\end{tikzpicture}
}
  \caption{Dual-engine decision flowchart.
  Mode \texttt{auto} prefers the compile path; warm disk/process \ir{} avoids
  repeated \jit{} cost; nopython kernels execute when available, otherwise
  object-mode; hard failures fall back to interpretation.}
  \label{fig:dual}
\end{figure}

\subsection{Bar-loop data flow}
\label{sec:dataflow}

Figure~\ref{fig:barloop} details data movement inside a run.
\ohlcv{} columns are bound as builtins; the host iterates $\baridx$; script
evaluation updates the series store, invokes incremental or full \ta{} kernels,
and drives the strategy engine.
Outputs are packed into plots, drawings, and events.

\begin{figure}[t]
  \centering
  \resizebox{0.98\textwidth}{!}{% Bar-loop data-flow: OHLCV → series → plots / strategy events
\begin{tikzpicture}[
  font=\sffamily\footnotesize,
  box/.style={draw, rounded corners=2pt, align=center, inner sep=5pt, minimum height=0.95cm},
  store/.style={box, fill=pyneblue!10, draw=pyneblue!70!black, minimum width=2.3cm},
  procstep/.style={box, fill=pyneorange!12, draw=pyneorange!70!black, minimum width=2.5cm},
  outbox/.style={box, fill=pynegreen!12, draw=pynegreen!70!black, minimum width=2.2cm},
  arr/.style={-{Stealth[length=2mm]}, thick, pynegray},
  node distance=0.45cm and 0.55cm
]
  % Inputs
  \node[store] (o) {open};
  \node[store, right=of o] (h) {high};
  \node[store, right=of h] (l) {low};
  \node[store, right=of l] (c) {close};
  \node[store, right=of c] (v) {volume};
  \node[fit=(o)(h)(l)(c)(v), draw=pyneblue!35, dashed, inner sep=5pt,
        label={[font=\sffamily\tiny\color{pynegray}]above:\ohlcv{} columns (length $n$)}] (ohlcv) {};

  \node[procstep, below=0.95cm of ohlcv] (bind) {Bind builtins\\(\texttt{close}, \texttt{bar\_index}, \ldots)};
  \node[procstep, below=of bind] (loop) {For $\baridx = 0 \ldots n-1$};
  \node[procstep, below=of loop, minimum width=8.5cm] (body) {%
    Evaluate script body once per bar\\
    series push / history $[k]$ / \texttt{var} retain / \texttt{varip} retain
  };

  \node[store, below left=0.85cm and 0.1cm of body] (ser) {Series store\\capped / ring};
  \node[procstep, below=0.85cm of body] (ta) {\ta{} kernels\\(inc.\ / full)};
  \node[procstep, below right=0.85cm and 0.1cm of body] (strat) {Strategy engine\\orders / fills};

  \node[outbox, below=0.9cm of ser] (plots) {plots\\fills\\bgcolor};
  \node[outbox, below=0.9cm of ta] (draw) {drawings\\lines / labels};
  \node[outbox, below=0.9cm of strat] (ev) {entries / exits\\alerts / equity};

  \draw[arr] (ohlcv) -- (bind);
  \draw[arr] (bind) -- (loop);
  \draw[arr] (loop) -- (body);
  \draw[arr] (body) -- (ser);
  \draw[arr] (body) -- (ta);
  \draw[arr] (body) -- (strat);
  \draw[arr] (ser) -- (plots);
  \draw[arr] (ta) -- (draw);
  \draw[arr] (strat) -- (ev);
  \draw[arr] (ta.west) -- (ser.east);
  \draw[arr] (ta.east) -- (strat.west);

  \node[draw, rounded corners, fill=codebg, below=0.55cm of draw, minimum width=9.2cm, align=left, font=\ttfamily\tiny] (note) {%
    na-propagation: arithmetic/compare with na $\Rightarrow$ na;\\
    short-circuit and/or; request.security: same-symbol simple OHLCV else na
  };
\end{tikzpicture}
}
  \caption{Bar-loop data flow from \ohlcv{} inputs through series storage and
  \ta{}/strategy engines to plots, drawings, and events.
  The bottom annotation summarizes $\na{}$ and \code{request.security}
  policies.}
  \label{fig:barloop}
\end{figure}

Algorithm~\ref{alg:barloop} captures the interpret host in pseudocode.
The compile path replaces the body of the inner evaluation with kernel calls
but preserves the same outer envelope construction.

\begin{algorithm}[t]
\caption{Interpret bar-loop host (simplified)}
\label{alg:barloop}
\begin{algorithmic}[1]
\Require script AST $P$, bars $B[0..n-1]$, options $\Omega$
\Ensure response envelope $\mathcal{E}$
\State $\textit{store} \gets \mathrm{InitSeriesStore}(\Omega)$
\State $\textit{strat} \gets \mathrm{InitStrategy}(\Omega)$
\State $\textit{plots} \gets \emptyset$;\ \ $\textit{events} \gets \emptyset$
\For{$\baridx \gets 0$ \textbf{to} $n-1$}
  \State $\mathrm{BindBuiltins}(\textit{store}, B[\baridx], \baridx)$
  \State $\mathrm{BeginBar}(\textit{strat}, B[\baridx])$
  \State $\mathrm{EvalBody}(P, \textit{store}, \textit{strat}, \textit{plots})$
  \State $\mathrm{ProcessPendingFills}(\textit{strat}, B[\baridx])$
  \State $\mathrm{EndBar}(\textit{strat}, \textit{events})$
  \State $\mathrm{CapSeries}(\textit{store})$ \Comment{optional memory bound}
\EndFor
\State $\mathcal{E} \gets \mathrm{PackEnvelope}(\textit{plots}, \textit{events}, \textit{meta})$
\State \textbf{return} $\mathcal{E}$
\end{algorithmic}
\end{algorithm}

\subsection{Warm compile and IR cache}
\label{sec:warm}

Cold \numba{} compilation can dominate latency for short interactive runs.
\pyne{} therefore implements:

\begin{itemize}[leftmargin=1.2em]
  \item \textbf{Disk \ir{} cache} keyed by compilation inputs (source digest,
    option flags, library versions as applicable).
  \item \textbf{Process prewarm} via \cli{}/\api{} hooks that force
    compilation of common kernels at worker start.
  \item \textbf{Corrupt cache recovery} that discards invalid entries and
    recomputes rather than failing the request path.
\end{itemize}

Internal Round~7 measurements report warm compile on the order of
$\sim 0.01\,\mathrm{ms}$ for a representative \texttt{ta\_combo} artifact, with
run latency $\sim 1.06\,\mathrm{ms}$ at $5{,}000$ bars (Section~\ref{sec:eval}).

\subsection{Incremental technical analysis kernels}
\label{sec:inc-ta}

Many indicators admit $O(1)$ or amortized updates per bar given prior
state~\citep{baker2012incremental,wilder1978new,bollinger2001,kama1998kaufman}.
Both engines exploit incremental forms for families such as
\code{sma}/\code{ema}/\code{rma}, \code{rsi}, \code{macd}, \code{bb},
\code{kama}, \code{stoch}, and others.
Kernel microbenchmarks in Round~7 show large per-kernel speedups (e.g.,
\texttt{kama} $\sim 121\times$, \texttt{bb} $\sim 217\times$ versus naive
recompute styles in the measured configuration).
Numerical agreement is checked against full recompute within floating-point
tolerance~\citep{higham2002accuracy,press2007numerical}.

\subsection{Strategy engine}
\label{sec:strategy}

Strategy support includes entries and exits, commission and slippage models,
one-cancels-other style constraints where implemented, and pending-fill
behaviour under pyramiding limits.
A representative semantic choice: when pyramiding is disabled or capped,
pending fills may average according to documented VWAP-style rules rather than
silently dropping intents.
Risk helpers and order-fill tests form part of the regression suite.

\subsection{Security policy for \texttt{request.security}}
\label{sec:security-policy}

Multi-symbol and multi-timeframe requests are a major complexity axis of the
language~\citep{tradingview-pine}.
\pyne{} adopts a conservative default:

\begin{definition}[Security resolution policy]
\label{def:security}
For \code{request.security} (and related forms), \emph{same-symbol simple
\ohlcv{}} expressions may resolve against the bound data feed.
Foreign symbols or complex expressions that would require inventing unavailable
market state resolve to $\na$ rather than fabricating closes.
\end{definition}

This policy prioritizes honesty for offline and CI contexts over incomplete
emulation of a full multi-asset platform data graph.

\subsection{Response envelope and parity harness}
\label{sec:envelope}

Regardless of engine, a successful run returns a structured envelope:

\begin{itemize}[leftmargin=1.2em]
  \item \textbf{Plots / fills / bgcolor} series aligned to bar indices.
  \item \textbf{Drawings} with garbage-collection limits
    (\code{max\_lines\_count}, \code{max\_labels\_count}, \ldots).
  \item \textbf{Strategy events} (entries, exits, equity samples).
  \item \textbf{Alerts} from \code{alert}/\code{alertcondition} with frequency
    semantics (\code{once\_per\_bar}, \code{once\_per\_bar\_close}, \code{all}).
  \item \textbf{Metadata} (engine used, timings, warnings, mode fallback flags).
\end{itemize}

The \emph{parity harness} compares interpret vs.\ compile outputs on shared
fixtures, asserting series alignment within absolute/relative tolerances
$(\abserr, \relerr)$ and matching event shapes.
Parity here means \emph{internal engine consistency}, not certification against
a proprietary platform.

%==============================================================================
\section{Product Surfaces}
\label{sec:surfaces}
%==============================================================================

\subsection{CLI}
\label{sec:cli}

The \code{pynescript} \cli{} supports check, format, lint, compile, run, data
fetch, and prewarm subcommands.
It is the primary interface for local scripts and automation:

\begin{lstlisting}[style=python,caption={Programmatic evaluation sketch.},label={lst:eval-api}]
# Conceptual host usage (simplified)
# runtime.run(source, ohlcv, mode="auto") -> envelope
\end{lstlisting}

Container images publish multi-arch \cli{} builds for reproducible desks and
CI.

\subsection{Language server}
\label{sec:lsp-surface}

\texttt{pynescript-lsp} implements \lsp{}~3.x features~\citep{lsp2023}:
diagnostics (parse + lint), completion, hover signatures, go-to-definition,
references, document symbols, semantic tokens, and formatting via unparse.
Distribution options include pip extras (\texttt{hoox-pyne[lsp]}) and Nuitka
onefile binaries~\citep{nuitka} bundled by the VS~Code extension.

\subsection{Pro API}
\label{sec:pro-api}

The Flask Pro backend~\citep{flask} exposes authenticated HTTP endpoints,
notably:

\begin{itemize}[leftmargin=1.2em]
  \item \texttt{POST /run} --- evaluate script with data and options;
  \item \texttt{POST /preview/chart} --- render chart artifacts;
  \item \texttt{POST /backtest/quick} --- strategy summary paths.
\end{itemize}

Middleware enforces \api{} keys and tier limits.
Workers may prewarm compile caches to meet interactive latency SLOs.

\subsection{Editors and edge}
\label{sec:edge}

\begin{itemize}[leftmargin=1.2em]
  \item \textbf{VS~Code extension:} first-class \texttt{.pyne}/\texttt{.pine}
    associations and bundled \lsp{} binary.
  \item \textbf{Neovim / Zed / Emacs:} client configs under \file{clients/}.
  \item \textbf{pine-worker / pyne-worker:} TypeScript and Python edge hosts
    sharing evaluate contracts~\citep{cloudflare-workers}.
  \item \textbf{Pyodide / AXIS:} browser-side evaluation for documentation and
    lightweight apps~\citep{pyodide}.
\end{itemize}

\subsection{Surface comparison}
\label{sec:surface-table}

Table~\ref{tab:surfaces} compares major product surfaces along packaging,
primary users, and evaluation capability.

\begin{table}[t]
\centering
\caption{Product surface comparison for \pyne{} and related hosts.}
\label{tab:surfaces}
\begin{tabular}{@{}llp{4.6cm}@{}}
\toprule
\textbf{Surface} & \textbf{Package / binary} & \textbf{Role} \\
\midrule
Core library & \texttt{hoox-pyne} / \texttt{pynescript} & Parse, \ast{}, evaluate, compile \\
\cli{} & \texttt{pynescript} & Desk automation, CI, prewarm \\
\lsp{} & \texttt{pynescript-lsp} & Editor intelligence ($\sim$800+ builtins) \\
VS~Code ext. & VSIX & Bundled \lsp{}, language config \\
Pro \api{} & Flask / containers & HTTP run, chart, backtest \\
pine-worker & Bun / TS & Edge evaluate (partial builtins) \\
pyne-worker & Python worker & Thin edge host over core semantics \\
Pyodide/AXIS & browser & In-page demos and docs \\
\bottomrule
\end{tabular}
\end{table}

%==============================================================================
\section{Evaluation and Deployment Experience}
\label{sec:eval}
%==============================================================================

Unless noted, performance figures are \emph{internal measurements} collected
on development hardware during 2026 performance campaigns (Rounds~1--7 and
related agent runs).
They are representative engineering benchmarks, not standardized public
leaderboards.

\subsection{Microbenchmark methodology}
\label{sec:method}

Scripts include minimal empty indicators, single-indicator programs
(\code{ta.sma}, \code{ta.rsi}, \code{ta.macd}), multi-indicator combinations,
and strategy-shaped fixtures.
Bar counts of $2{,}000$ and $5{,}000$ are common.
We report median wall times over repeated runs after optional warm-up.
Compile paths distinguish cold \jit{} vs.\ warm cache hits.

\subsection{Front-end performance}
\label{sec:perf-fe}

Early agent campaigns measured:

\begin{itemize}[leftmargin=1.2em]
  \item \textbf{Parse:} approximately $5.4\times$ improvement on a complex
    $\sim$16\,KB script via SLL-first parsing with LL fallback, cheaper
    annotation handling, and faster location attachment.
  \item \textbf{Unparse:} approximately $1.95\times$ on a multi-script corpus
    via thread-local unparser reuse and type-keyed dispatch.
\end{itemize}

An \ast{} LRU cache keyed by source digest further reduces multi-parse latency
in \lsp{} and repeated \api{} submissions (multi-parse speedups up to three
orders of magnitude on cache hits in Round~7 notes).

\subsection{Interpret versus compile}
\label{sec:perf-engines}

For numeric \ta{} scripts, compile+execute can be
$\sim 49$--$110\times$ faster than interpret in measured MACD/multi
configurations after incremental kernels and plot packing improvements.
Interpret-side incremental \ta{} yields roughly $1.9$--$6.5\times$ gains on
relevant Runtime paths versus non-incremental baselines.

Table~\ref{tab:bench} and Figure~\ref{fig:bench} summarize representative
median latencies.
Interpret rows reflect bar-loop host costs at a few thousand bars; compile
rows reflect warm nopython execution style latencies (run phase).
Exact values vary by CPU, Numba version, and plot export settings; the table
is intended to communicate magnitude and ranking.

\begin{table}[t]
\centering
\caption{Representative median latencies (internal measurements, 2026).
Interpret figures are bar-loop style at $\sim$2k bars for combo; compile run
figures are warm kernels at $\sim$5k bars where noted.}
\label{tab:bench}
\begin{tabular}{@{}lrrrl@{}}
\toprule
\textbf{Workload} & \textbf{Interpret (ms)} & \textbf{Compile run (ms)} & \textbf{Approx.\ ratio} & \textbf{Notes} \\
\midrule
minimal @2k & 15.2 & --- & --- & Round~7 interpret \\
\texttt{ta\_sma} @2k & 23.6 & 0.04 & $\sim$590$\times$ & compile @5k SMA \\
\texttt{ta\_rsi} & $\sim$22 & 0.08 & large & warm compile \\
\texttt{ta\_macd} & $\sim$35 & 0.12 & large & warm compile \\
MULTI & $\sim$48 & 0.21 & $\sim$230$\times$ & multi-indicator \\
\texttt{ta\_combo} & 152.9 & 1.06 & $\sim$144$\times$ & interp@2k / comp@5k \\
strategy\_ish @2k & 68.9 & (object mix) & --- & fills + events \\
\midrule
warm compile only & --- & $\sim$0.01 & --- & cache hit, combo \\
\bottomrule
\end{tabular}
\end{table}

\begin{figure}[t]
  \centering
  \resizebox{\textwidth}{!}{% Benchmark bar chart: interpret vs compile latency (representative internal measurements, 2026)
\begin{tikzpicture}
\begin{axis}[
  ybar,
  bar width=9pt,
  width=0.92\textwidth,
  height=6.2cm,
  ylabel={Median latency (ms)},
  symbolic x coords={SMA, RSI, MACD, MULTI, ta\_combo},
  xtick=data,
  ymin=0,
  ymax=180,
  legend style={at={(0.5,1.02)}, anchor=south, legend columns=2, draw=none, font=\footnotesize},
  nodes near coords,
  nodes near coords align={vertical},
  every node near coord/.append style={font=\tiny, /pgf/number format/fixed, /pgf/number format/precision=2},
  enlarge x limits=0.18,
  grid=major,
  grid style={dashed, gray!30},
  tick label style={font=\small},
  label style={font=\small},
  axis background/.style={fill=codebg},
]
\addplot[fill=pyneblue!70, draw=pyneblue!90!black] coordinates {
  (SMA, 18.5) (RSI, 22.0) (MACD, 35.0) (MULTI, 48.0) (ta\_combo, 152.85)
};
\addplot[fill=pyneorange!80, draw=pyneorange!90!black] coordinates {
  (SMA, 0.04) (RSI, 0.08) (MACD, 0.12) (MULTI, 0.21) (ta\_combo, 1.06)
};
\legend{Interpret (cold path style), Compile+exec (warm)}
\end{axis}
\end{tikzpicture}
}
  \caption{Interpret vs.\ warm compile+exec median latency for representative
  numeric workloads (internal measurements, 2026).
  Axes are not normalized to identical bar counts across all series; the chart
  emphasizes order-of-magnitude gaps for kernel-heavy scripts.}
  \label{fig:bench}
\end{figure}

\subsection{Memory: series cap and ring buffer}
\label{sec:memory}

Long historical runs accumulate series history.
Enabling a series cap (default on in recent builds) reduced long-run list
memory by approximately $78\times$ in Round~7 experiments relative to uncapped
growth.
An optional chronological ring buffer further bounds residency; default rollout
is gated on corpus greenness.

\subsection{Compatibility posture}
\label{sec:compat}

Table~\ref{tab:compat} summarizes compatibility themes across the Python core
and edge hosts.
The core \texttt{pynescript} implementation is the reference for parity tests;
workers may delegate or reimplement subsets.

\begin{table}[t]
\centering
\caption{Compatibility matrix summary (qualitative).
Full feature tables live in project documentation~\citep{pynescript2026}.}
\label{tab:compat}
\begin{tabular}{@{}p{3.2cm}ccc@{}}
\toprule
\textbf{Area} & \textbf{pynescript} & \textbf{pine-worker} & \textbf{pyne-worker} \\
\midrule
Parser (v5/v6) & Full (ANTLR) & Full (generated TS) & Delegated \\
\ast{} coverage & Full ASDL & Partial (Zod) & Delegated \\
\na{} / series & Full & Full (subset builtins) & Delegated \\
Strategy fills & Full & Full (evolving) & Full (via core) \\
Compile engine & Numba dual path & N/A & N/A \\
\code{request.security} & Same-symbol simple; else $\na$ & Policy-dependent & Delegated \\
Parity vs core & Reference & Partial fixtures & Fixture-oriented \\
\bottomrule
\end{tabular}
\end{table}

Corpus parse rates on sanitized public-style script sets have been measured in
the high nineties of percent for supported syntax after lexer/sanitize
hardening; residual failures concentrate on exotic scrapes and unsupported
platform-only constructs.

\subsection{Test suite scale}
\label{sec:tests}

The repository maintains a substantial \texttt{pytest} suite spanning parser
round-trips, linter cases, \ta{} goldens, strategy fills, compiler kernels,
interpret$\leftrightarrow$compile parity, \lsp{} features, and backend routes.
Recent verification notes cite on the order of $600+$ passing core tests after
Round~7 merges, with additional edge-worker suites elsewhere.
Tests prefer deterministic synthetic \ohlcv{} over live network data.

\subsection{Deployment lessons}
\label{sec:deploy}

Operational experience suggests:

\begin{enumerate}[leftmargin=1.4em]
  \item Prefer \textbf{warm compile workers} on the Pro path; interpret remains
    valuable for debugging and for scripts outside the nopython surface.
  \item Treat \textbf{plot packing} as a first-class cost center: profiles show
    export/registry work dominating some multi-plot interpret runs even when
    kernels are cheap.
  \item Keep \textbf{parse caches} coherent with call-site metadata scrubbing so
    that cache hits cannot poison empty plot registries.
  \item Document \textbf{security $\na$ policy} prominently: users expecting
    full multi-symbol platform behaviour must supply data explicitly.
\end{enumerate}

%==============================================================================
\section{Related Work}
\label{sec:related}
%==============================================================================

\subsection{DSL toolchains and compiler construction}
\label{sec:rw-dsl}

Domain-specific languages are a long-standing software engineering
topic~\citep{hudak1996dsl,fowler2010dsl,erwig2011dsl}.
\pyne{} follows the classical pipeline of lexing, parsing, \ast{}
construction, and multi-backend execution~\citep{aho2006dragon,cooper2011engineering,muchnick1997advanced}.
Using \antlr{}~\citep{parr2013antlr} and \asdl{}~\citep{wang1997asdl} mirrors
choices in other language implementations that separate declarative syntax
from semantic passes.
Visitor-based lowering resembles patterns in production compilers and teaching
interpreters~\citep{gamma1994design,nystrom2021crafting}.

\subsection{Financial and statistical time-series languages}
\label{sec:rw-fin}

Technical analysis packages in general-purpose languages (e.g., NumPy-centric
stacks~\citep{harris2020numpy,brooks2015python}) provide library APIs rather
than a charting \textsc{Dsl}.
\pine{}~\citep{tradingview-pine} is distinctive in coupling series indexing,
plot effects, and strategy simulation in one language.
\pyne{} reimplements that \emph{language-shaped} interface offline without
claiming platform equivalence.
Classical indicator definitions remain the numerical
foundation~\citep{wilder1978new,bollinger2001,murphy1999technical,kama1998kaufman}.

\subsection{Dual interpreter + JIT systems}
\label{sec:rw-jit}

Many production language VMs combine interpretation with selective
compilation~\citep{chambers1991iterative,peytonjones2003ghc,lattner2004llvm}.
\pyne{}'s dual engine is narrower: a domain bar-loop host with \numba{}
kernels~\citep{lam2015numba} for numeric subgraphs and an \ast{} interpreter
for full language coverage.
The design goal is pragmatic performance for indicators, not a general tracing
\jit{} for arbitrary Python.

\subsection{Language servers and editor ecosystems}
\label{sec:rw-lsp}

The Language Server Protocol~\citep{lsp2023,voutilainen2016lsp} decouples
editors from language implementations.
\pyne{}'s pygls server~\citep{murphy2020pygls} follows this model, enabling
VS~Code, Neovim, Zed, and Emacs from one codebase---an explicit alternative to
browser-only authoring.

\subsection{Edge and browser runtimes}
\label{sec:rw-edge}

Deploying evaluation to Cloudflare Workers~\citep{cloudflare-workers} and
Pyodide~\citep{pyodide} continues a trend of moving compute closer to clients.
\pyne{}'s contribution is a shared evaluate contract across these hosts rather
than a single-runtime monoculture.

%==============================================================================
\section{Discussion and Limitations}
\label{sec:discussion}
%==============================================================================

\subsection{No full platform parity}
\label{sec:lim-parity}

\pyne{} does not reproduce the entire \tv{} product: proprietary charting UI,
closed data services, social features, and some platform-only builtins are out
of scope.
Numerical results on community scripts can differ when hosts implement
different fill models, session calendars, or security resolution.
Users should treat \pyne{} as an open evaluation and tooling substrate, not as
a drop-in legal or functional substitute for the commercial platform.

\subsection{Foreign \texttt{request.security} policy}
\label{sec:lim-security}

Resolving foreign or complex \code{request.security} calls to $\na$
(Definition~\ref{def:security}) avoids silent wrong answers but surprises users
who expect multi-asset charts without providing data.
Future work includes richer multi-feed binders under explicit user control.

\subsection{Compile surface incompleteness}
\label{sec:lim-compile}

Not every language feature lowers to nopython.
Object-mode and interpret fallbacks preserve correctness at the cost of
latency cliffs when a script leaves the fast path.
Continuing kernel coverage and strategy lowering remains active engineering.

\subsection{AGPL implications}
\label{sec:lim-agpl}

\pyne{} is licensed under \textsc{Agpl}-3.0-or-later~\citep{agpl3}.
Network-facing services that modify and provide the software to users inherit
corresponding source-offer obligations.
Organizations embedding \pyne{} in SaaS should review compliance with counsel;
the license choice intentionally favours open reciprocity for a community
toolchain~\citep{hoox2026}.

\subsection{Trademark and independence}
\label{sec:lim-tm}

\pine{} and \tv{} are trademarks of TradingView, Inc.
\pyne{} is independent and unofficial.
Documentation references public language materials for interoperability only
and does not redistribute proprietary platform software.

\subsection{Threats to validity for benchmarks}
\label{sec:lim-bench}

Reported speedups depend on hardware, library versions, plot export settings,
and script mix.
Median times from internal campaigns should be read as engineering evidence of
relative engine behaviour, not as guaranteed SLAs.

%==============================================================================
\section{Conclusion}
\label{sec:conclusion}
%==============================================================================

We presented \pyne{}, an open dual-engine toolchain for offline \pine{}
evaluation.
By structuring the system as a classical language pipeline---\antlr{} front-end,
\asdl{} \ast{}, bar-loop host with interpret and \numba{} compile modes---and
by sharing that core across \cli{}, \lsp{}, Pro \api{}, and edge clients, the
project enables reproducible research and editor-native workflows outside
browser-locked environments.
Internal measurements show substantial gains from incremental kernels and warm
compilation for numeric scripts, while a conservative security policy and
parity harness emphasize honest offline semantics.
Limitations remain around full platform parity and compile coverage; we view
these as a roadmap rather than a claim of completion.
\pyne{} is available as \texttt{hoox-pyne} on PyPI under
\textsc{Agpl}-3.0-or-later as part of the \textsc{Hoox} open trading
stack~\citep{pynescript2026,hoox2026}.

%==============================================================================
\section*{Acknowledgments}
%==============================================================================

We thank contributors to the \textsc{Hoox}/\pyne{} ecosystem and users who
filed compatibility reports.
Performance campaigns (Rounds~1--7) benefited from automated multi-agent
benchmarking harnesses and regression corpora maintained in-tree.

%==============================================================================
\bibliographystyle{plainnat}
\bibliography{../common/bibliography}

\end{document}
